\chapter{Judgments of Learning: Delay, Cues and Predictive Accuracy}
\label{sec:p19-judgments-of-learning}

A judgment of learning predicts later memory performance for recently
studied material. Its accuracy depends on what is predicted, when the
judgment is made, what information is displayed, and how accuracy is
scored. Being able to rank two items by eventual recall is different
from assigning the right numerical probability to either item.

This chapter uses Nelson and Dunlosky's 1991 delayed-judgment paper
\cite{NelsonDunlosky1991DelayedJOL} and Dunlosky and Nelson's 1992
cue comparison \cite{DunloskyNelson1992CueJOL}. Author order changes
between the papers. Their experiments use paired unrelated nouns and
university students, not classroom vocabulary instruction.
The evidence-record conventions of
Chapter~\ref{sec:p19-evidence-methods} and the calibration distinctions
of Chapter~\ref{sec:p19-calibration-extreme-confidence} apply.

\section{Monitoring, performance and the score}

\begin{hbdefinition}{def:p19-jol-record}{An item-level monitoring record}
For one participant in one specified condition, let \(n\) be a positive
integer and index the studied items by \(i=1,\ldots,n\).
Record a judgment \(q_i\in[0,1]\) for each item and an eventual
criterion outcome \(y_i\in\{0,1\}\). The judgment predicts success
under the stated future test, cue and timing. Define performance
as \(\bar y=n^{-1}\sum_i y_i\), mean judgment as
\(\bar q=n^{-1}\sum_iq_i\), and mean judgment bias as
\(\bar q-\bar y\). Positive bias means overprediction on this record.
Missing judgments and unscored tests require a stated inclusion rule.

For these papers, a JOL is a prospective confidence rating, whereas
the criterion is later cued recall of the response word. Monitoring
is the judgment process; control is a subsequent decision such as
which item to restudy. The experiments' fixed study presentations
do not themselves test every proposed self-paced control policy.
\end{hbdefinition}

\begin{hbdefinition}{def:p19-jol-gamma}{Concordance, ties and Goodman--Kruskal gamma}
Consider unordered pairs \(i<j\). A pair is concordant if
\((q_i-q_j)(y_i-y_j)>0\), discordant if this product is negative,
and excluded from this comparison if the product is zero.
Write \(N_+\) and \(N_-\) for the concordant and discordant counts.
Thus a tie in the judgment, outcome, or both contributes to neither.
When \(N_++N_->0\), define
\[
 G=\frac{N_+-N_-}{N_++N_-}.
\]
When the denominator is zero, \(G\) is undefined. Assigning it zero
would be a new scoring convention, not the value of the fraction.
The selected papers compute gamma separately for each person's
immediate and delayed items, then summarize those person-level values.
\end{hbdefinition}

\begin{hbtheorem}{thm:p19-jol-gamma}{proved=annex:p19.jol-gamma}{What gamma does and does not express}
Defined gamma lies in \([-1,1]\). If a pair is sampled uniformly
from the pairs with different judgments and different outcomes,
the conditional probability that the recalled item has the higher
judgment is
\[
 \frac{N_+}{N_++N_-}=\frac{1+G}{2}.
\]
For binary outcomes, gamma is undefined exactly when every outcome
is the same or every judgment is the same.
A function strictly increasing on the set of judgment values preserves
gamma when applied to every judgment. Even when its output lies in
\([0,1]\), it need not preserve calibration, mean bias or squared error.
\end{hbtheorem}

The probability in this theorem is conditional on two exclusions:
different judgments and different outcomes. It is not the probability
that a single item will be recalled, the proportion of all pairs
ordered correctly, or the percentage of outcome variance explained.
Gamma equal to one need not imply calibrated probability forecasts.
Changing the number or pattern of ties changes which pairs enter
the denominator. A claim about gamma must therefore retain the
rating scale, outcome variation and undefined-score handling.

The word resolution also needs qualification. Here within-person
relative monitoring accuracy concerns ordering items by their future
outcomes. The calibration chapter's
\(\operatorname{Res}=\mathbb E[(m(C)-\mu)^2]\) is a different
quantity under a declared joint law. Neither is simply overall recall.

\section{The 1991 main experiment}
\label{sec:p19-jol-1991-main}

Thirty University of Washington undergraduates volunteered for extra
course credit. They studied 66 unrelated concrete noun pairs using an
Apple II computer, at eight seconds per pair. The first six were
practice items and were not tested. The remaining sixty formed two
blocks of thirty. Within each block, fifteen items were randomly
assigned to immediate JOL and fifteen to delayed JOL, with no more
than three consecutive items assigned to the same condition.
This was a within-person item comparison, not thirty independently
assigned groups \cite[p.~268]{NelsonDunlosky1991DelayedJOL}.

The JOL displayed the stimulus word alone and asked for the predicted
likelihood of recalling its response when prompted later by that
stimulus, approximately ten minutes hence. The displayed confidence
levels were 0, 20, 40, 60, 80 and 100 percent. Judgments and recall were
self-paced and typed. An immediate JOL followed the pair's offset.
Delayed judgments occurred after the block's study and immediate
judgments: the first five delayed items were judged in random order,
then the second five, then the final five. Study and JOL were
separated by 10--33 other study or judgment events.
This manipulation is not a fixed, separately randomized
thirty-second interval.

After the last delayed JOL in block two, recall began, with block one
tested before block two and item order randomized again within blocks.
The paper reports no intervening ten-minute filler in this main
procedure. The query's approximate future time is not a statement
that every item had an identical JOL-to-test interval.

Here SEM denotes the reported standard error of the mean.
Mean recall was 45.5 percent, SEM 5.1 percentage points, for immediate
items and 45.9 percent, SEM 5.3, for delayed items.
Mean gamma was \(+.38\), SEM .08, for immediate JOLs and \(+.90\),
SEM .02, for delayed JOLs; the paired direction comparison used a
sign test with \(p<.001\).
Five participants had an indeterminate gamma in one condition and
were excluded from this relative-accuracy analysis, leaving 25.
The reported condition SEMs do not supply a paired-difference SEM
without the within-person covariance
\cite[p.~268 and note~2]{NelsonDunlosky1991DelayedJOL}.

Figure~1's delayed distribution had no gamma below \(+.63\);
more than 60 percent attained at least \(+.90\), and 80 percent
at least \(+.80\). Every delayed gamma exceeded the immediate
condition's mean, which is different from claiming that every
participant's paired difference was positive.
The mean delayed gamma maps to .95 under the conditional-pair
interpretation; this does not mean 95 percent of studied items
were recalled \cite[p.~269]{NelsonDunlosky1991DelayedJOL}.

For absolute accuracy, the authors grouped items receiving the
same predicted percentage and compared actual recall with that
percentage. Their Figure~2 places the delayed calibration curve
closer to the diagonal. The fitted descriptive lines, with both
axes in percentage units, were
\[
 y=.84x+10\quad\text{(delayed)},\qquad
 y=.39x+28\quad\text{(immediate)}.
\]
These are fits to calibration-curve points, not participant-level
causal regressions. The article does not provide all bin denominators,
raw records, or confidence intervals for the curve or line parameters.
Closeness of a fitted line is not a proof of exact calibration at
every rating \cite[p.~269, note~3]{NelsonDunlosky1991DelayedJOL}.

\section{The 1991 interval checks and reported extensions}
\label{sec:p19-jol-1991-controls}

One alternative explanation was that delayed judgments simply
occurred nearer the recall test. Within the main experiment,
the authors compared delayed items in block one with immediate
items in block two. More intervening block-two study followed the
former judgments, yet their mean gammas were \(+.90\) and \(+.36\),
respectively, with \(p<.001\) by a sign test.

The continued second footnote reports a separate experiment with
twenty new participants. Its procedure was otherwise the same,
but recall tested block two before block one. In the critical
comparison, delayed block-one JOLs consequently had a longer
JOL-to-recall interval than immediate block-two JOLs.
Recall was 29 versus 40 percent, \(t(19)=2.78,p<.05\).
Mean gammas were \(+.78\) versus \(+.47\), with medians \(+.98\)
versus \(+.54\); the delayed advantage had \(p<.05\) by a sign test.
Thus greater monitoring association coexisted with lower recall
in this comparison. The note does not print condition SEMs or
an undefined-gamma exclusion count for this second experiment
\cite[pp.~268--269, note~2]{NelsonDunlosky1991DelayedJOL}.

Retrospective end-session interviews in that twenty-person
experiment found nineteen reporting attempted recall during delayed
JOLs, versus half reporting that strategy during immediate JOLs.
These are strategy reports, not randomized assignments to retrieval
processes or direct observations of every judgment.
The discussion also mentions other replications, including median
gammas \(+1.0\) versus \(+.43\) when delayed JOLs had six intervening
items. It gives no complete method or sample size for that brief
report; it is not treated here as a fully described additional trial
\cite[p.~270 and note~4]{NelsonDunlosky1991DelayedJOL}.

\section{The 1992 cue comparison}
\label{sec:p19-jol-1992-design}

Dunlosky and Nelson recruited one hundred University of Washington
undergraduates for extra credit, tested individually.
The paper describes random assignment, by order of appearance,
of fifty to stimulus-alone JOL cues and fifty to
stimulus--response-pair cues. It does not print the assignment
sequence or randomization algorithm.
Immediate versus delayed JOL remained a within-person factor.
Thus cue type varied between people and JOL timing within people
\cite[p.~375]{DunloskyNelson1992CueJOL}.

The material again comprised 66 unrelated concrete noun pairs,
with six untested practice items and two blocks of thirty.
Write \(C_{\rm norm}\) for the word-norm concreteness score
denoted \(C\) in the source. Its reported selection criterion
was \(C_{\rm norm}\geq6.08\). The underlying norm construction
and complete item list are not supplied by these two articles;
6.08 is a selection threshold, not a mean JOL or a recall rate
\cite[p.~375]{DunloskyNelson1992CueJOL}.
Study now lasted six seconds per pair. Each block had fifteen
items assigned to each timing condition, with the same
no-more-than-three-consecutive constraint. Delayed items were
judged by successive thirds of their block, with at least ten
other study/JOL events intervening. The rating scale and intended
future cued-recall question used the same six printed percentage levels.

After all judgments, an unrelated filler task lasted ten minutes,
then recall proceeded in successive groups of ten items from
study order, randomized within each group. Participants had to
type a response; omissions were not allowed. A response was
scored correct when its first three letters were correct, to
reduce the influence of spelling. This criterion is not silently
replaced by exact whole-word spelling. The added filler, six-second
presentation and recall ordering distinguish this procedure from
the 1991 main experiment.

\subsection{Relative monitoring accuracy}

Five stimulus-alone participants and six pair-cue participants
had indeterminate gammas and were excluded from relative-accuracy
analyses. The remaining analysis groups were 45 and 44.
The reported values below are medians, unlike the principal
1991 mean-gamma comparison:
\[
\begin{array}{l|cc}
 \text{JOL cue}&\text{Immediate }G&\text{Delayed }G\\\hline
 \text{Stimulus alone}&+.45&+.93\\
 \text{Stimulus--response pair}&+.55&+.60
\end{array}
\]
Every one of the 45 defined-score stimulus-alone participants had
a larger delayed gamma; the sign test gave \(p<.001\).
The pair-cue sign test gave \(p=.18\).
That latter result does not establish a zero timing effect or
statistical equivalence.

Delayed stimulus-alone gamma also exceeded immediate pair-cue
gamma, with reported Mann--Whitney \(U=1464,p<.001\), and delayed
pair-cue gamma, with \(U=1424,p<.001\).
The text accompanying Figure~1 reports that more than 60 percent of the delayed
stimulus-alone group reached \(G\geq.90\), compared with fewer
than 20 percent in each other condition.
The figure plots four cumulative distributions of person-level
gammas, at thresholds \(+.90,+.80,+.70,+.60,+.50\), followed
by the remaining below-\(.50\) category that brings each curve
to 100 percent. Its vertical coordinate is the percentage of
people attaining a level or higher, not the percentage of items
recalled and not the proportion of explained variance.
No interval estimates for these gamma contrasts are printed.
The medians' difference is not asserted to equal the median
of the individual paired differences
\cite[p.~376]{DunloskyNelson1992CueJOL}.

\subsection{Calibration and rating use}

For each condition, the authors first grouped a participant's
items by predicted percentage, computed recall for that group,
and then plotted mean actual recall across participants.
They describe delayed stimulus-alone judgments as better
calibrated than the other three conditions. The four fitted
lines, again in percentage units, were
\[
\begin{array}{l|l}
 \text{Condition}&\text{Fitted actual recall}\\\hline
 \text{Delayed, stimulus alone}&7.4+.79x\\
 \text{Immediate, stimulus alone}&17.9+.40x\\
 \text{Delayed, pair cue}&22.4+.58x\\
 \text{Immediate, pair cue}&17.0+.53x
\end{array}
\]
Figure~2 compares immediate and delayed calibration within each
cue type: stimulus alone in the upper panel and the pair cue
in the lower panel. Figure~3 reuses the delayed stimulus-alone
curve as the reference against immediate pair cues in its upper
panel and delayed pair cues in its lower panel. These are
comparisons of the same four conditions, not additional samples.
Both axes run in percentage units, with points at the six
reported judgment levels and the diagonal as the calibration
benchmark. All four displayed conditions show some recall at
zero-percent predictions and less than complete recall at
hundred-percent predictions. The plotted improvement therefore
does not establish exact calibration or smaller error at every
individual rating \cite[pp.~376--377, figs.~1--3]{DunloskyNelson1992CueJOL}.

Notes~3--4 supply the fitted-line parameters above
\cite[p.~380]{DunloskyNelson1992CueJOL}.
The paper does not supply all participant-by-rating denominators
or a complete missing-bin convention needed to reconstruct the
curves and their uncertainty.

The Table~1 investigation addresses whether near-perfect gamma
merely reflected use of two ratings. The complete printed table
contains the following five rating-count columns.
The final column below is their calculated sum:
\[
\begin{array}{l|rrrrr|r}
 &\multicolumn{5}{c|}{\text{Number of different ratings used}}&
 \text{Listed total}\\
 \text{Condition}&1&2&3&4&5&\\\hline
 \text{Stimulus alone, delayed}&0&4&13&13&8&38\\
 \text{Stimulus alone, immediate}&0&1&8&17&14&40\\
 \text{Pair cue, delayed}&1&7&6&13&14&41\\
 \text{Pair cue, immediate}&0&1&8&20&14&43
\end{array}
\]
The source itself prints no sixth column, although the rating
scale permits six distinct ratings. The listed totals are not
the fifty recruited participants per cue group; nor are they
the 45 and 44 participants in the defined-gamma analyses.
These printed counts alone do not identify the unlisted
participants or their rating-count categories.
In particular, the first row prints 34 people in categories
three through five, while the accompanying prose reports that
46 of 50 used at least three ratings.
The prose claim is retained as reported, but no unprinted
six-rating cell is supplied as an observed count.

The paper reports the four-condition comparison as
\(\chi^2(15)=20.45,p>.10\).
Failure to reject that comparison does not prove identical
rating distributions or exclude every influence of ties and
rating use. Immediate and delayed rows also repeat the same
people. The incomplete printed frequency information and that
repeated-person structure do not supply a complete reconstruction
of the reported analysis or a participant-level equivalence test
\cite[pp.~377--378, table~1]{DunloskyNelson1992CueJOL}.

\subsection{Recall was a different outcome}

For all fifty participants in each cue group, the source reports
the following mean recall proportions and SEMs:
\[
\begin{array}{l|cc}
 \text{Cue}&\text{Immediate}&\text{Delayed}\\\hline
 \text{Stimulus alone}&.33\ (.03)&.31\ (.03)\\
 \text{Stimulus--response pair}&.32\ (.03)&.44\ (.04)
\end{array}
\]
The overall mean was .35; the displayed condition range is
.31--.44. In the reported analysis of variance,
\(\mathrm{MS}_e\) denotes the error mean square.
The timing main effect was
\(F(1,98)=9.72,\mathrm{MS}_e=.013,p<.01\);
the cue main effect was
\(F(1,98)=1.88,\mathrm{MS}_e=.094,p=.17\);
and the timing-by-cue interaction was
\(F(1,98)=23.61,\mathrm{MS}_e=.013,p<.001\).

For pair cues the delayed recall advantage had
\(t(49)=5.49,p<.001\). For stimulus-alone cues, the immediate
versus delayed difference had \(t(49)=1.27,p=.21\).
Delayed pair-cue recall also exceeded immediate stimulus-alone
recall, \(t(98)=2.26,p<.05\), and delayed stimulus-alone recall,
\(t(98)=2.74,p<.01\).
Thus the condition with the strongest reported monitoring
association was not the condition with the highest recall.
The gamma exclusions do not reduce these fifty-person recall
comparisons to 45 or 44 participants
\cite[p.~378]{DunloskyNelson1992CueJOL}.

\section{A fully constructed monitoring calculation}
\label{sec:p19-jol-worked}

\begin{hbexample}{def:p19-jol-worked}{Six fictional items and their actual denominators}
The following record is invented solely to illustrate the scores:
\[
\begin{array}{c|rrrrrr}
i&1&2&3&4&5&6\\\hline
q_i&0&.2&.4&.4&.8&1\\
y_i&0&1&0&1&1&1
\end{array}
\]
There are four recalled and two unrecalled items, hence eight
cross-outcome pairs. The pair of items 3 and 4 has tied judgments.
The remaining seven consist of six concordant pairs and one
discordant pair, giving \(G=5/7\) and conditional concordance
probability \(6/7\). Neither denominator is six items or fifteen
unrestricted unordered pairs.

Recall is \(2/3\), mean prediction \(7/15\), and mean bias
\(-1/5\). Under the equal-item empirical law, the exact judgment
classes have observed success rates \(0,1,1/2,1,1\) at
ratings \(0,.2,.4,.8,1\), respectively.
Their squared calibration error is \(7/60\) and the single-coordinate
squared forecast loss is \(1/5\).
The good positive association does not remove those calibration errors.
\end{hbexample}

\begin{hbtheorem}{thm:p19-jol-calculations}{proved=annex:p19.jol-calculations}{The worked scores and two comparison boundaries}
The calculations in the six-item example hold.
For outcomes \((0,0,1,1)\), changing perfectly calibrated reports
\((0,0,1,1)\) to \((.99,.99,1,1)\) preserves \(G=1\), but gives
mean overprediction .495 and squared calibration error .49005.

If two fictional people report \((.8,.9)\) and \((.1,.2)\),
respectively, each with outcomes \((0,1)\), both within-person
gammas equal one. Pooling their four items gives gamma \(1/2\).
The pooled mean judgment equals mean recall, but the squared
calibration error is \(13/40\).
Thus pooling item pairs, averaging person-level gammas and
assessing calibration are different operations.
\end{hbtheorem}

\section{What these experiments support}
\label{sec:p19-jol-inference}

The timing and cue manipulations provide comparisons under their
specified procedures in the recruited students. They do not make
the volunteers representative of all learners. Within-person
item assignment does not turn all sixty responses into independent
participants. The relative-accuracy comparisons additionally condition
on gamma being defined, with exclusions that may depend on the
manipulated cue and resulting response distributions.
Complete raw records and confidence intervals are not supplied.

The authors propose monitoring of short- and long-term memory as
one explanation: immediate accessibility may be uninformative about
later recall, and displaying the response again may recreate that
problem. They also propose similarity between monitoring and test
contexts. The 1991 proposal used an approximately thirty-second
short-term-information timescale; the experiments did not independently
establish that value as a universal threshold.
Neither account is uniquely identified by the two experiments.
The 1992 discussion predicts a different cue advantage for recognition
of complete pairs, as distinct from recognizing either word alone,
and proposes progressively revealed response letters as a hybrid cue.
Those are proposed tests, not implemented conditions
in this corpus \cite[pp.~379--380]{DunloskyNelson1992CueJOL}.

Cue exposure can also change the later criterion. Showing a response
again can support additional learning; attempting retrieval can affect
later access. The 1992 authors discuss maintenance of already available
responses, restoration of unavailable responses, and possible extra
learning of weak items as distinct accounts.
Their final note also records a proposed contrast between judgments
based on subjective experience and theory-based judgments.
These suggestions are candidates for further tests, not additional
mechanisms isolated experimentally in this corpus
\cite[p.~380, note~6]{DunloskyNelson1992CueJOL}.
Their second note mentions another unreported experiment consistent
with the pair-cue null result, including a nonreliable trend favoring
immediate pair-cue judgments. It gives no sample, effect estimate
or full procedure. It cannot be treated as a fully described
independent replication.

The fully described experiments have no no-JOL control condition, and neither study directly
randomizes the proposed latent retrieval or memory-store mechanism.
The observed relation concerns recall after the actual cue/judgment
procedure, not a counterfactual memory state unaffected by that procedure.
The retrospective strategy reports do not close that identification gap.

For practical use, state whether the aim is ranking items for possible
restudy, forecasting numerical recall probabilities, or increasing
retention. Preserve the criterion, cue, delay and rating scale when
transferring an evidence claim. The articles' suggestions about foreign
language learning and more efficient self-paced study motivate further
tests; their noun-pair results do not establish a universal optimum
delay, a universal learning prescription, or an improvement in every
one of these outcomes simultaneously.
